Production LLM agents in 2025 and 2026 spend most of their compute
in a tool loop. Untrusted text from emails, retrieval hits, web
pages, Model Context Protocol servers, and other agents gets
concatenated into the context window. That context then decides
which tool runs next, with which arguments. This is fragile in a
specific, well-studied way: indirect prompt injection
\citep{greshake2023, liu2023prompt}. An attacker who controls a
single document in the agent's retrieval-augmented context can
subvert privileged tools through instruction smuggling, and no
amount of careful prompt engineering at the host side fixes it.

The standard defense, and the one almost every OSS guardrail library
has chosen, is to treat injection as a classification problem.
Llama Prompt Guard 2, ShieldGemma, Llama Guard 4
\citep{inan2023llamaguard}, OpenAI Moderation, Azure AI Content
Safety, and dozens of community classifiers sit somewhere on this
spectrum. But the AgentDojo evaluation
\citep{debenedetti2024agentdojo} shows the ceiling: even with a
classifier in front, GPT-4o drops from 69\% utility to 45\% under
injection, and attack success rates on canonical patterns like
``important-message'' stay above 50\%. The structural problem is
that an attacker with adversarial input can iterate against a
classifier, and the classifier has bounded capacity to tell content
apart from instruction.

The research literature, meanwhile, has moved toward a different
abstraction. CaMeL \citep{camel2025} treats indirect injection as
an information-flow problem. Untrusted data carries labels
(provenance, capability, integrity), and the runtime enforces that
no untrusted-derived control flow can reach a privileged tool.
Where classifier defenses degrade gracefully under attack, IFC
defeats indirect injection by construction: a tool whose arguments
carry an UNTRUSTED label simply does not execute. CaMeL reports
near-zero attack success on AgentDojo at roughly 2.7$\times$ token
cost.

None of this work ships as a library you can install with one line.
That alone would be enough motivation for AIsafePy.flow. But IFC
isn't the only gap. Two others matter equally in 2026.

\paragraph{Streaming guardrails are still wrong by default.}
Voice agents, IDE assistants, and agentic CLIs stream tokens.
Existing guardrails either buffer the full response and rerun a
classifier (sacrificing time-to-first-token) or run guards in
parallel and risk leaking violative tokens before cancellation.
A streaming-native design that forecasts likely continuations and
stops at the first signal of harm is the right shape, but no
open-source library packages this end-to-end with a tiered cascade,
explicit latency budgets, and white-box activation probes for
self-hosted models.

\paragraph{Eval-to-guardrail compilation is a manual workflow.}
PyRIT, Garak, and Inspect AI produce JSONL logs of failures.
Promoting those failures into runtime guards (classifiers, regexes,
policy rules, steering vectors) is currently a two-engineer-day
process. Galileo Luna-2 and Lakera Red$\rightarrow$Guard sell this
automation as a commercial moat. No open-source equivalent exists.

\paragraph{Contributions.} This paper presents AIsafePy, an
Apache-2.0 Python library that targets all three gaps. The
contributions are:

\begin{enumerate}[leftmargin=2em, itemsep=0.25em]
  \item The first OSS implementation of CaMeL-style capability-based
    IFC as drop-in adapters for OpenAI Agents SDK, LangGraph,
    LlamaIndex, Anthropic tools, and MCP servers
    (Section~\ref{sec:design}, \texttt{aisafepy.flow}).

  \item A three-tier streaming guardrail pipeline with an explicit
    p95 latency budget, structured \texttt{GuardDecision} outputs,
    and optional white-box activation probes for self-hosted models
    (Section~\ref{sec:design}, \texttt{aisafepy.stream}).

  \item A continuous eval-to-guardrail compiler that emits four
    runtime artifact modalities (classifier, regex, policy, steering
    vector) from PyRIT, Garak, Inspect, and Langfuse sources
    (Section~\ref{sec:design}, \texttt{aisafepy.adapt}).

  \item An MCP-aware tool-poisoning defense combining static
    description screening with hashed-manifest drift detection,
    aligned with the Model Context Protocol specification
    \citep{mcpspec} and community-reported attack patterns
    against MCP servers.

  \item An evaluation roadmap and a public benchmarking harness
    targeting AgentDojo, HarmBench, and a streaming-guard latency
    benchmark, with reproducibility instructions in the repository.
\end{enumerate}

Section~\ref{sec:related} positions AIsafePy against prior work.
Section~\ref{sec:design} describes the system design.
Section~\ref{sec:implementation} covers implementation choices.
Section~\ref{sec:evaluation} presents the evaluation roadmap.
Section~\ref{sec:discussion} lists what the library does not solve.
